% ---------------------------------------------------------------------
%  CHƯƠNG 2 — CƠ SỞ LÝ THUYẾT VÀ CÔNG NGHỆ
% ---------------------------------------------------------------------
\chapter{Cơ sở lý thuyết và công nghệ}
\label{ch:coso}

\begin{muctieu}
Chương này trình bày lý thuyết về hệ gợi ý và cách tính độ tương đồng dùng
trong đồ án, cùng danh mục công cụ và thư viện để xây dựng hệ thống.
\end{muctieu}

\section{Kiến thức nền tảng}

\begin{dinhnghia}[Hệ gợi ý]
Hệ gợi ý \en{(recommender system)} là hệ thống dự đoán mức độ phù hợp giữa
một người dùng $u$ và một đối tượng $i$ trong tập đối tượng $\mathcal{I}$,
từ đó xếp hạng và đề xuất những đối tượng có điểm phù hợp cao nhất cho
người dùng.
\end{dinhnghia}

Hai hướng tiếp cận phổ biến nhất là lọc dựa trên nội dung
\en{(content-based filtering)}, dựa vào đặc trưng của chính đối tượng và hồ
sơ nhu cầu của người dùng, và lọc cộng tác \en{(collaborative filtering)},
dựa vào hành vi của những người dùng có sở thích tương tự. Với bài toán tìm
phòng trọ, số lượt tương tác của mỗi người dùng thường rất ít vì nhu cầu
thuê trọ không lặp lại thường xuyên, nên lọc cộng tác thuần tuý dễ gặp vấn
đề khởi động lạnh \en{(cold start)}. Đồ án chọn lọc dựa trên nội dung làm
thành phần chính, vì mỗi phòng trọ luôn có đầy đủ đặc trưng mô tả ngay khi
được đăng, không phụ thuộc vào lịch sử tương tác.

Mỗi phòng trọ $i$ được biểu diễn bằng một vec-tơ đặc trưng $\mathbf{x}_i$
gồm các thành phần đã chuẩn hoá: giá thuê, khu vực, diện tích, danh sách
tiện ích và mô tả văn bản vec-tơ hoá theo TF-IDF. Nhu cầu của người dùng
cũng được biểu diễn bằng một vec-tơ $\mathbf{q}$ cùng cấu trúc, ghép từ
ngân sách mong muốn, khu vực ưu tiên, diện tích tối thiểu và các tiện ích
bắt buộc. Độ phù hợp giữa nhu cầu và một phòng trọ được tính bằng tổng có
trọng số của độ tương đồng cô-sin trên từng nhóm đặc trưng
\begin{equation}
\mathrm{score}(\mathbf{q}, \mathbf{x}_i) =
  w_1 \cdot \mathrm{sim}_{\text{giá}} +
  w_2 \cdot \mathrm{sim}_{\text{khu vực}} +
  w_3 \cdot \mathrm{sim}_{\text{tiện ích}} +
  w_4 \cdot \mathrm{sim}_{\text{mô tả}},
\label{eq:score}
\end{equation}
trong đó mỗi $\mathrm{sim}(\cdot)$ là độ tương đồng cô-sin
\begin{equation}
\mathrm{sim}(\mathbf{a}, \mathbf{b}) =
  \frac{\mathbf{a} \cdot \mathbf{b}}{\lVert \mathbf{a} \rVert \, \lVert \mathbf{b} \rVert}
\label{eq:cosine}
\end{equation}
tính trên nhóm thành phần tương ứng, và $w_1, \dots, w_4$ là trọng số không
âm có tổng bằng một. Các phòng trọ được xếp hạng giảm dần theo điểm số ở
Công thức~\eqref{eq:score} và trả về $k$ phòng có điểm cao nhất.

\begin{ghinho}
Việc chuẩn hoá từng nhóm đặc trưng về cùng thang giá trị trước khi tính độ
tương đồng cô-sin là bắt buộc, nếu không thành phần có giá trị tuyệt đối
lớn hơn, ví dụ giá thuê tính bằng đồng, sẽ lấn át hoàn toàn các thành phần
còn lại trong Công thức~\eqref{eq:score}.
\end{ghinho}

\section{Công cụ và thư viện}

\begin{table}[H]
\centering
\caption{Công cụ và thư viện sử dụng trong đồ án.}
\label{tab:congcu}
\begin{tabular*}{\textwidth}{@{\extracolsep{\fill}}>{\bfseries\color{daunavy}}l >{\color{daumaroon}}l l}
\toprule
\rowcolor{gray100}\bfseries\color{daunavy}Thành phần & \bfseries\color{daumaroon}Phiên bản & \bfseries\color{daunavy}Vai trò \\
\midrule
Python       & 3.11 & Ngôn ngữ lập trình chính \\
pandas       & 2.2  & Đọc và xử lý dữ liệu phòng trọ dạng bảng \\
scikit-learn & 1.4  & Vec-tơ hoá TF-IDF và tính độ tương đồng cô-sin \\
Flask        & 3.0  & Xây dựng giao diện web và API gợi ý \\
SQLite       & 3.x  & Lưu trữ dữ liệu phòng trọ và người dùng \\
Leaflet.js   & 1.9  & Hiển thị bản đồ vị trí phòng trọ trên nền OpenStreetMap \\
\bottomrule
\end{tabular*}
\end{table}

\section{Các công trình liên quan}

Các hệ thống gợi ý bất động sản và phòng cho thuê thường đi theo hai hướng
chính. Hướng thứ nhất dùng lọc dựa trên nội dung với đặc trưng thủ công như
giá, vị trí và diện tích, có ưu điểm dễ triển khai và không cần dữ liệu
tương tác lớn~\cite{RussellNorvig2021}. Hướng thứ hai kết hợp thêm lọc cộng
tác để tận dụng hành vi của những người dùng có sở thích tương tự, giúp
tăng tính đa dạng của kết quả gợi ý nhưng đòi hỏi lượng dữ liệu tương tác đủ
lớn để tránh khởi động lạnh~\cite{HoangKiem}. Đồ án này chọn hướng thứ nhất
làm nền tảng, đồng thời bổ sung một thành phần cộng tác đơn giản dựa trên
những phòng trọ mà người dùng có hồ sơ nhu cầu tương tự đã lưu lại, nhằm
giảm bớt hạn chế của lọc dựa trên nội dung thuần tuý là kết quả gợi ý có xu
hướng quá giống với tiêu chí đã nhập, thiếu bất ngờ hữu ích.
